
\documentclass{article}
\usepackage{colortbl}
\usepackage{makecell}
\usepackage{multirow}
\usepackage{supertabular}

\begin{document}

\newcounter{utterance}

\centering \large Interaction Transcript for game `cladder', experiment `full\_v1.5\_default', episode 9064 with llama{-}3.1{-}8b{-}instruct.
\vspace{24pt}

{ \footnotesize  \setcounter{utterance}{1}
\setlength{\tabcolsep}{0pt}
\begin{supertabular}{c@{$\;$}|p{.15\linewidth}@{}p{.15\linewidth}p{.15\linewidth}p{.15\linewidth}p{.15\linewidth}p{.15\linewidth}}
    \# & \multicolumn{2}{c}{Player} && \multicolumn{2}{c}{Game Master} \\
    \hline

    \theutterance \stepcounter{utterance}  
    & & & \multicolumn{4}{p{0.6\linewidth}}{
        \cellcolor[rgb]{0.9,0.9,0.9}{
            \makecell[{{p{\linewidth}}}]{
                \texttt{\tiny{[P1$\langle$GM]}}
                \texttt{You are an expert in causal inference. The following question is not a typical commonsense query, but rather a meticulously designed question created by a professor specializing in causal inference, intended to assess the students' mastery of the course content.} \\
\texttt{Imagine a self{-}contained, hypothetical world with only the following conditions, and without any unmentioned factors or causal relationships: Pexu has a direct effect on hwax. Hwax has a direct effect on rukz. Method 1: We look at how pexu correlates with rukz case by case according to hwax. Method 2: We look directly at how pexu correlates with rukz in general. To understand how pexu affects rukz, is it more correct to use the Method 1 than Method 2?} \\
\texttt{Start your answer with 'Yes' or 'No', followed by additional reasoning or evidence to support your explanation.} \\
            }
        }
    }
     \\ \\

    \theutterance \stepcounter{utterance}  
    & \multicolumn{4}{p{0.6\linewidth}}{
        \cellcolor[rgb]{0.561,0.737,0.561}{
            \makecell[{{p{\linewidth}}}]{
                \texttt{\tiny{[P1$\rangle$GM]}}
                \texttt{No, it is not more correct to use Method 1 than Method 2 to understand how pexu affects rukz.} \\
\\ 
\texttt{In this hypothetical world, we have a direct effect of pexu on hwax and a direct effect of hwax on rukz. This creates a causal chain: pexu → hwax → rukz. When we look directly at the correlation between pexu and rukz (Method 2), we are observing the total effect of pexu on rukz, which includes both the direct effect of pexu on rukz (if any) and the indirect effect through hwax.} \\
\\ 
\texttt{On the other hand, Method 1 involves looking at the correlation between pexu and rukz case by case according to hwax. This is essentially a form of stratification or subgroup analysis, where we are conditioning on the value of hwax to examine the relationship between pexu and rukz. However, this approach can lead to confounding bias, as the observed correlation between pexu and rukz within each stratum of hwax may not accurately reflect the underlying causal relationship.} \\
\\ 
\texttt{In particular, if there is a non{-}linear relationship between hwax and rukz, or if there are interactions between pexu and hwax that affect rukz, then the stratified analysis (Method 1) may not provide a consistent or unbiased estimate of the causal effect of pexu on rukz. In contrast, the direct analysis (Method 2) can provide a more comprehensive and unbiased understanding of the total effect of pexu on rukz, including both the direct and indirect effects through hwax.} \\
\\ 
\texttt{Therefore, in this hypothetical world, it is generally more correct to use Method 2 than Method 1 to understand how pexu affects rukz.} \\
            }
        }
    }
    & & \\ \\

    \theutterance \stepcounter{utterance}  
    & & & \multicolumn{2}{p{0.3\linewidth}}{
        \cellcolor[rgb]{0.9,0.9,0.9}{
            \makecell[{{p{\linewidth}}}]{
                \texttt{\tiny{[GM$|$GM]}}
                \texttt{No} \\
            }
        }
    }
    & & \\ \\

    \theutterance \stepcounter{utterance}  
    & & & \multicolumn{2}{p{0.3\linewidth}}{
        \cellcolor[rgb]{0.9,0.9,0.9}{
            \makecell[{{p{\linewidth}}}]{
                \texttt{\tiny{[GM$|$GM]}}
                \texttt{no} \\
            }
        }
    }
    & & \\ \\

    \theutterance \stepcounter{utterance}  
    & & & \multicolumn{2}{p{0.3\linewidth}}{
        \cellcolor[rgb]{0.9,0.9,0.9}{
            \makecell[{{p{\linewidth}}}]{
                \texttt{\tiny{[GM$|$GM]}}
                \texttt{game\_result = WIN} \\
            }
        }
    }
    & & \\ \\

\end{supertabular}
}

\end{document}
